\documentclass[12pt,a4paper]{report}

% Packages
\usepackage[utf8]{inputenc}
\usepackage[T1]{fontenc}
\usepackage{geometry}
\usepackage{graphicx}
\usepackage{listings}
\usepackage{xcolor}
\usepackage{hyperref}
\usepackage{fancyhdr}
\usepackage{titlesec}
\usepackage{tocloft}
\usepackage{float}
\usepackage{tabularx}
\usepackage{booktabs}
\usepackage{enumitem}
\usepackage{amsmath}
\usepackage{caption}
\usepackage{subcaption}

% Page geometry
\geometry{margin=2.5cm}

% Colors
\definecolor{codegreen}{rgb}{0,0.6,0}
\definecolor{codegray}{rgb}{0.5,0.5,0.5}
\definecolor{codepurple}{rgb}{0.58,0,0.82}
\definecolor{backcolour}{rgb}{0.95,0.95,0.92}
\definecolor{darkblue}{rgb}{0.0,0.0,0.6}

% Code listing style
\lstdefinestyle{mystyle}{
    backgroundcolor=\color{backcolour},   
    commentstyle=\color{codegreen},
    keywordstyle=\color{magenta},
    numberstyle=\tiny\color{codegray},
    stringstyle=\color{codepurple},
    basicstyle=\ttfamily\footnotesize,
    breakatwhitespace=false,         
    breaklines=true,                 
    captionpos=b,                    
    keepspaces=true,                 
    numbers=left,                    
    numbersep=5pt,                  
    showspaces=false,                
    showstringspaces=false,
    showtabs=false,                  
    tabsize=2,
    frame=single
}
\lstset{style=mystyle}

% Hyperref setup
\hypersetup{
    colorlinks=true,
    linkcolor=darkblue,
    filecolor=magenta,      
    urlcolor=blue,
    pdftitle={Multi-Threaded Text Analyzer - Project Report},
    pdfauthor={Merouane Bensadallah},
}

% Header and footer
\pagestyle{fancy}
\fancyhf{}
\fancyhead[L]{\leftmark}
\fancyhead[R]{Multi-Threaded Text Analyzer}
\fancyfoot[C]{\thepage}
\renewcommand{\headrulewidth}{0.4pt}
\renewcommand{\footrulewidth}{0.4pt}

% Title formatting
\titleformat{\chapter}[display]
{\normalfont\huge\bfseries}{\chaptertitlename\ \thechapter}{20pt}{\Huge}

% Document info
\title{
    \vspace{-2cm}
    \includegraphics[width=0.3\textwidth]{logo_placeholder.png}\\[1cm]
    \textbf{\Huge Multi-Threaded Text Analyzer}\\[0.5cm]
    \Large Desktop Application Project Report\\[1cm]
    \large Concurrency Programming Course (2025-2026)
}
\author{
    \textbf{Merouane Bensadallah}\\
    \textbf{Bensalem Hadil}\\
    \textbf{Zerizer Mellisa}\\[0.3cm]
    \texttt{GitHub: MarwaneQ}\\[0.5cm]
    \textbf{Instructor:} Dr. Boudjellal Nada\\
    \textbf{Institution:} NTIC - GL
}
\date{\today}

\begin{document}

% Title page
\begin{titlepage}
\centering
\vspace*{2cm}

{\Huge\bfseries Multi-Threaded Text Analyzer\par}
\vspace{0.5cm}
{\Large Desktop Application Project Report\par}
\vspace{1.5cm}

{\large Concurrency Programming Course (2025-2026)\par}
\vspace{2cm}

\begin{tabular}{rl}
\textbf{Students:} & Merouane Bensadallah \\
                   & Bensalem Hadil \\
                   & Zerizer Mellisa \\
\textbf{GitHub:} & MarwaneQ \\
\textbf{Instructor:} & Dr. Boudjellal Nada \\
\textbf{Institution:} & NTIC - GL \\
\end{tabular}

\vspace{2cm}

{\large \today\par}

\vfill

\begin{abstract}
This report documents the development of a Multi-Threaded Text Analyzer desktop application built with Java 21 and JavaFX 21. The project demonstrates concurrent programming concepts through the implementation of a file analysis system capable of processing multiple text files simultaneously. The development was organized into three sprints: UI Development, Single-Threaded Analysis, and Multi-Threaded Analysis with ThreadPoolExecutor.
\end{abstract}

\end{titlepage}

% Table of contents
\tableofcontents
\newpage

% List of figures
\listoffigures
\newpage

%========================================
% CHAPTER 1: INTRODUCTION
%========================================
\chapter{Introduction}

\section{Project Overview}
The Multi-Threaded Text Analyzer is a JavaFX-based desktop application designed to analyze multiple text files concurrently. The application provides comprehensive text statistics including word count, character count, unique words, and word frequency analysis.

\section{Objectives}
The main objectives of this project are:
\begin{itemize}
    \item Develop a user-friendly graphical interface for text file analysis
    \item Implement single-threaded text analysis algorithms
    \item Introduce multi-threading for concurrent file processing
    \item Ensure thread safety and proper resource management
    \item Provide real-time progress tracking and performance metrics
\end{itemize}

\section{Technology Stack}
\begin{table}[H]
\centering
\begin{tabularx}{\textwidth}{|l|X|l|}
\hline
\textbf{Component} & \textbf{Technology} & \textbf{Version} \\
\hline
Programming Language & Java (LTS) & 21 \\
\hline
UI Framework & JavaFX & 21 \\
\hline
Build Tool & Apache Maven & 3.11.0 \\
\hline
Module System & Java Platform Module System (JPMS) & Java 21 \\
\hline
\end{tabularx}
\caption{Technology Stack}
\end{table}

\section{Project Structure}
\begin{lstlisting}[language=bash, caption=Project Directory Structure]
text-analyzer/
├── src/main/java/com/textanalyzer/
│   ├── App.java                     # Application entry point
│   ├── module-info.java             # JPMS module configuration
│   ├── controllers/                 # UI controllers (MVC)
│   │   ├── PrimaryController.java   # Main screen logic
│   │   └── SecondaryController.java # Results screen logic
│   ├── services/                    # Business logic layer
│   │   ├── TextAnalysisService.java # Analysis algorithms
│   │   ├── MultiFileAnalysisService.java # Concurrent analysis
│   │   └── ThreadPoolManager.java   # Thread management
│   ├── models/                      # Data models (POJOs)
│   │   ├── AnalysisResult.java      # Single file results
│   │   └── MultiFileAnalysisResult.java # Multi-file results
│   ├── utils/                       # Helper utilities
│   └── exceptions/                  # Custom exceptions
│       └── InvalidFileException.java
├── src/main/resources/com/textanalyzer/
│   ├── primary.fxml                 # Main UI layout
│   ├── secondary.fxml               # Results UI layout
│   └── styles.css                   # CSS styling
└── pom.xml                          # Maven configuration
\end{lstlisting}

%========================================
% CHAPTER 2: SPRINT 1 - UI DEVELOPMENT
%========================================
\chapter{Sprint 1: User Interface Development}

\section{Sprint Objectives}
The first sprint focused on creating a functional and user-friendly GUI skeleton:
\begin{enumerate}
    \item Design the user interface using JavaFX
    \item Implement file loading functionality
    \item Display selected files in the application
    \item Add event handlers for UI interactions
\end{enumerate}

\section{UI Design}

\subsection{Design Philosophy}
The UI was designed with a modern aesthetic featuring:
\begin{itemize}
    \item \textbf{Glassmorphism effects} - Semi-transparent cards with blur
    \item \textbf{Gradient backgrounds} - Purple to pink color scheme
    \item \textbf{Smooth animations} - Hover and click effects
    \item \textbf{Responsive layout} - BorderPane with flexible regions
\end{itemize}

\subsection{Color Scheme}
\begin{table}[H]
\centering
\begin{tabularx}{\textwidth}{|l|X|}
\hline
\textbf{Element} & \textbf{Color/Gradient} \\
\hline
Header & \texttt{linear-gradient(135deg, \#667eea → \#764ba2)} \\
\hline
Background & Light gray (\texttt{\#f5f5f5}) \\
\hline
Cards & White with soft shadows \\
\hline
Load Button & Blue gradient (\texttt{\#4a90e2 → \#357abd}) \\
\hline
Start Button & Green gradient (\texttt{\#11998e → \#38ef7d}) \\
\hline
\end{tabularx}
\caption{UI Color Scheme}
\end{table}

\section{Implementation Details}

\subsection{Primary View (Main Screen)}
The main screen contains:
\begin{itemize}
    \item Header with application title
    \item File list with custom cell renderer
    \item Load Files button with file picker
    \item Start Analysis button
    \item Status bar with indicators
\end{itemize}

\subsection{Custom File List Cell}
\begin{lstlisting}[language=Java, caption=Custom ListCell Implementation]
private class FileListCell extends ListCell<String> {
    private final HBox container;
    private final Label fileNameLabel;
    private final Button removeButton;
    
    public FileListCell() {
        container = new HBox(10);
        container.setAlignment(Pos.CENTER_LEFT);
        
        fileNameLabel = new Label();
        fileNameLabel.setMaxWidth(Double.MAX_VALUE);
        HBox.setHgrow(fileNameLabel, Priority.ALWAYS);
        
        removeButton = new Button("x");
        removeButton.getStyleClass().add("remove-button");
        removeButton.setOnAction(e -> {
            String fileName = getItem();
            if (fileName != null) {
                removeFile(fileName);
            }
        });
        
        container.getChildren().addAll(fileNameLabel, removeButton);
    }
    
    @Override
    protected void updateItem(String fileName, boolean empty) {
        super.updateItem(fileName, empty);
        if (empty || fileName == null) {
            setGraphic(null);
        } else {
            fileNameLabel.setText(fileName);
            setGraphic(container);
        }
    }
}
\end{lstlisting}

\subsection{File Loading}
\begin{lstlisting}[language=Java, caption=File Loading with FileChooser]
@FXML
public void onLoadFiles() {
    FileChooser fileChooser = new FileChooser();
    fileChooser.setTitle("Select Text Files to Analyze");
    fileChooser.getExtensionFilters().add(
        new FileChooser.ExtensionFilter("Text Files", "*.txt")
    );

    List<File> files = fileChooser.showOpenMultipleDialog(null);

    if (files != null && !files.isEmpty()) {
        fileListView.getItems().clear();
        fileNameToPathMap.clear();
        
        for (File file : files) {
            String fileName = file.getName();
            fileListView.getItems().add(fileName);
            fileNameToPathMap.put(fileName, file.getAbsolutePath());
        }

        updateStatus("Files loaded successfully", "status-ready");
        analyzeBtn.setDisable(false);
    }
}
\end{lstlisting}

\section{Sprint 1 Deliverables}
\begin{itemize}
    \item[$\checkmark$] Modern UI with glassmorphism design
    \item[$\checkmark$] File management with native file picker
    \item[$\checkmark$] Visual file list with individual remove buttons
    \item[$\checkmark$] Real-time status updates
    \item[$\checkmark$] Responsive design with minimum window constraints
    \item[$\checkmark$] CSS-based styling
\end{itemize}

%========================================
% CHAPTER 3: SPRINT 2 - SINGLE-THREADED
%========================================
\chapter{Sprint 2: Single-Threaded Analyzer}

\section{Sprint Objectives}
The second sprint focused on implementing text analysis functionality:
\begin{enumerate}
    \item Implement text analysis algorithms
    \item Display analysis results in the interface
    \item Handle errors gracefully
    \item Add progress bar visualization
    \item Ensure UI responsiveness during analysis
\end{enumerate}

\section{Analysis Algorithms}

\subsection{Word Counting}
\begin{lstlisting}[language=Java, caption=Word Counting Algorithm]
public int countWords(String content) {
    if (content == null || content.trim().isEmpty()) {
        return 0;
    }
    
    String[] words = content.trim().split("\\s+");
    int count = 0;
    for (String word : words) {
        String cleanWord = word.replaceAll("[^a-zA-Z0-9]", "");
        if (cleanWord.length() >= 3) {  // Filter short words
            count++;
        }
    }
    return count;
}
\end{lstlisting}

\subsection{Character Counting}
\begin{lstlisting}[language=Java, caption=Character Counting Algorithm]
public int countCharacters(String content) {
    if (content == null || content.trim().isEmpty()) {
        return 0;
    }
    // Remove all whitespace and count remaining characters
    return content.replaceAll("\\s+", "").length();
}
\end{lstlisting}

\subsection{Word Frequency Calculation}
\begin{lstlisting}[language=Java, caption=Word Frequency Algorithm]
public Map<String, Integer> calculateWordFrequencies(String content) {
    Map<String, Integer> frequencies = new HashMap<>();
    
    if (content == null || content.trim().isEmpty()) {
        return frequencies;
    }
    
    String[] words = content.toLowerCase().split("\\s+");
    
    for (String word : words) {
        String cleanWord = word.replaceAll("[^a-zA-Z0-9]", "");
        if (cleanWord.length() >= 3) {
            frequencies.put(cleanWord, 
                frequencies.getOrDefault(cleanWord, 0) + 1);
        }
    }
    
    return frequencies;
}
\end{lstlisting}

\subsection{Top N Frequencies}
\begin{lstlisting}[language=Java, caption=Top N Frequencies using Stream API]
private Map<String, Integer> getTopNFrequencies(
        Map<String, Integer> frequencies, int n) {
    return frequencies.entrySet()
        .stream()
        .sorted(Map.Entry.<String, Integer>comparingByValue().reversed())
        .limit(n)
        .collect(Collectors.toMap(
            Map.Entry::getKey,
            Map.Entry::getValue,
            (e1, e2) -> e1,
            LinkedHashMap::new
        ));
}
\end{lstlisting}

\section{JavaFX Task for Background Processing}
To keep the UI responsive during analysis, we use JavaFX's \texttt{Task} class:

\begin{lstlisting}[language=Java, caption=Background Task Implementation]
Task<AnalysisResult> analysisTask = new Task<>() {
    @Override
    protected AnalysisResult call() throws Exception {
        updateMessage("Reading file...");
        updateProgress(0.1, 1.0);
        
        AnalysisResult result = analysisService.analyzeFile(file);
        
        updateMessage("Complete!");
        updateProgress(1.0, 1.0);
        
        return result;
    }
    
    @Override
    protected void succeeded() {
        AnalysisResult result = getValue();
        SecondaryController.setAnalysisResult(result);
        App.setRoot("secondary");
    }
    
    @Override
    protected void failed() {
        Throwable exception = getException();
        SecondaryController.setError(exception.getMessage());
        App.setRoot("secondary");
    }
};

// Run in background thread
Thread analysisThread = new Thread(analysisTask);
analysisThread.setDaemon(true);
analysisThread.start();
\end{lstlisting}

\section{Threading Model - Sprint 2}
\begin{figure}[H]
\centering
\begin{verbatim}
┌─────────────────────────────────────────────────────────────┐
│                      MAIN UI THREAD                          │
│  (handles button clicks, animations, screen updates)         │
├─────────────────────────────────────────────────────────────┤
│  User clicks "Analyze" → Create Task → new Thread().start()  │
│         │                                    │               │
│         ▼                                    ▼               │
│  UI remains responsive              BACKGROUND THREAD        │
│  Progress bar animates              ┌────────────────────┐  │
│         │                           │ analyzeFile()      │  │
│         │                           │ countWords()       │  │
│         │                           │ calculateFreqs()   │  │
│         │◄────── onSucceeded() ◄────┴────────────────────┘  │
│         ▼                                                    │
│  Display results in UI                                       │
└─────────────────────────────────────────────────────────────┘
\end{verbatim}
\caption{Sprint 2 Threading Model}
\end{figure}

\section{Data Model}
\begin{lstlisting}[language=Java, caption=AnalysisResult Model]
public class AnalysisResult {
    private String fileName;
    private int totalWords;
    private int totalCharacters;
    private int uniqueWords;
    private List<String> uniqueWordsList;
    private Map<String, Integer> wordFrequencies;
    private long processingTimeMs;
    
    // Constructor, getters, setters...
}
\end{lstlisting}

\section{Sprint 2 Deliverables}
\begin{itemize}
    \item[$\checkmark$] Text analysis algorithms (word count, character count, frequencies)
    \item[$\checkmark$] Results display with statistics cards
    \item[$\checkmark$] Progress bar during analysis
    \item[$\checkmark$] Background thread execution with Task
    \item[$\checkmark$] Error handling with custom exceptions
    \item[$\checkmark$] Responsive UI during processing
\end{itemize}

%========================================
% CHAPTER 4: SPRINT 3 - MULTI-THREADED
%========================================
\chapter{Sprint 3: Multi-Threaded Analyzer}

\section{Sprint Objectives}
The third sprint introduced multi-threading for concurrent file analysis:
\begin{enumerate}
    \item Use Java concurrency tools (ExecutorService, ThreadPoolExecutor)
    \item Ensure thread safety and manage shared resources
    \item Test with different numbers of files
    \item Debug concurrency and performance issues
    \item Update UI for multi-file results
\end{enumerate}

\section{ThreadPoolManager Implementation}
\begin{lstlisting}[language=Java, caption=ThreadPoolManager Singleton]
public class ThreadPoolManager {
    private static ThreadPoolManager instance;
    private ExecutorService executorService;
    
    private static final int CORE_POOL_SIZE = 4;
    private static final int MAX_POOL_SIZE = 8;
    private static final long KEEP_ALIVE_TIME = 60L;
    
    private final AtomicInteger threadCounter = new AtomicInteger(0);
    
    private ThreadPoolManager() {
        initializeThreadPool();
    }
    
    public static synchronized ThreadPoolManager getInstance() {
        if (instance == null) {
            instance = new ThreadPoolManager();
        }
        return instance;
    }
    
    private void initializeThreadPool() {
        ThreadFactory threadFactory = runnable -> {
            Thread thread = new Thread(runnable);
            thread.setName("TextAnalyzer-Worker-" + 
                threadCounter.incrementAndGet());
            thread.setDaemon(true);
            return thread;
        };
        
        executorService = new ThreadPoolExecutor(
            CORE_POOL_SIZE,
            MAX_POOL_SIZE,
            KEEP_ALIVE_TIME,
            TimeUnit.SECONDS,
            new LinkedBlockingQueue<>(100),
            threadFactory,
            new ThreadPoolExecutor.CallerRunsPolicy()
        );
    }
    
    public <T> Future<T> submit(Callable<T> task) {
        return executorService.submit(task);
    }
}
\end{lstlisting}

\section{Multi-File Analysis Service}
\begin{lstlisting}[language=Java, caption=Concurrent File Analysis]
public MultiFileAnalysisResult analyzeFiles(List<File> files) {
    MultiFileAnalysisResult result = new MultiFileAnalysisResult();
    result.setTotalFilesProcessed(files.size());
    result.setOverallStartTime(System.currentTimeMillis());
    
    List<Future<AnalysisResult>> futures = new ArrayList<>();
    
    // Submit all files for concurrent processing
    for (File file : files) {
        Future<AnalysisResult> future = threadPoolManager.submit(() -> {
            String threadName = Thread.currentThread().getName();
            
            notifyProgress(result.getCompletedCount(), files.size(), 
                          file.getName(), "Analyzing...", threadName);
            
            try {
                AnalysisResult analysisResult = 
                    textAnalysisService.analyzeFile(file);
                return analysisResult;
            } catch (InvalidFileException e) {
                result.addFailure(file.getName(), e.getMessage());
                return null;
            }
        });
        futures.add(future);
    }
    
    // Collect results from all futures
    for (int i = 0; i < futures.size(); i++) {
        Future<AnalysisResult> future = futures.get(i);
        try {
            AnalysisResult analysisResult = 
                future.get(60, TimeUnit.SECONDS);
            if (analysisResult != null) {
                result.addResult(analysisResult);
            }
        } catch (TimeoutException e) {
            result.addFailure(files.get(i).getName(), "Timed out");
        } catch (Exception e) {
            result.addFailure(files.get(i).getName(), e.getMessage());
        }
    }
    
    result.setOverallEndTime(System.currentTimeMillis());
    return result;
}
\end{lstlisting}

\section{Thread-Safe Result Collection}
\begin{lstlisting}[language=Java, caption=Thread-Safe MultiFileAnalysisResult]
public class MultiFileAnalysisResult {
    private final List<AnalysisResult> successfulResults;
    private final Map<String, String> failedFiles;
    private final AtomicInteger completedCount;
    
    public MultiFileAnalysisResult() {
        this.successfulResults = new ArrayList<>();
        this.failedFiles = new ConcurrentHashMap<>();
        this.completedCount = new AtomicInteger(0);
    }
    
    // Thread-safe add method
    public synchronized void addResult(AnalysisResult result) {
        successfulResults.add(result);
        completedCount.incrementAndGet();
    }
    
    // Thread-safe failure recording
    public void addFailure(String fileName, String errorMessage) {
        failedFiles.put(fileName, errorMessage);
        completedCount.incrementAndGet();
    }
    
    // Speedup calculation
    public double getSpeedupFactor() {
        long parallel = getOverallProcessingTime();
        long sequential = getSequentialEquivalentTime();
        if (parallel == 0) return 0;
        return (double) sequential / parallel;
    }
}
\end{lstlisting}

\section{Multi-Threading Architecture}
\begin{figure}[H]
\centering
\begin{verbatim}
┌────────────────────────────────────────────────────────────────┐
│                    USER CLICKS "ANALYZE"                        │
└────────────────────────────────────────────────────────────────┘
                              │
                              ▼
┌────────────────────────────────────────────────────────────────┐
│   MAIN THREAD: Creates Task<MultiFileAnalysisResult>           │
│                Starts new Thread for the Task                   │
└────────────────────────────────────────────────────────────────┘
                              │
                              ▼
┌────────────────────────────────────────────────────────────────┐
│   BACKGROUND THREAD: MultiFileAnalysisService.analyzeFiles()   │
│                                                                 │
│   ┌──────────────────────────────────────────────────────┐    │
│   │         THREAD POOL (4-8 worker threads)              │    │
│   │                                                       │    │
│   │  ┌──────────┐  ┌──────────┐  ┌──────────┐           │    │
│   │  │ Worker-1 │  │ Worker-2 │  │ Worker-3 │  ...      │    │
│   │  │ file1.txt│  │ file2.txt│  │ file3.txt│           │    │
│   │  └──────────┘  └──────────┘  └──────────┘           │    │
│   └──────────────────────────────────────────────────────┘    │
│                                                                 │
│   Collects all Future<AnalysisResult> results                  │
└────────────────────────────────────────────────────────────────┘
                              │
                              ▼
┌────────────────────────────────────────────────────────────────┐
│   MAIN THREAD: onSucceeded() displays results in UI            │
│                Shows speedup factor, individual file results    │
└────────────────────────────────────────────────────────────────┘
\end{verbatim}
\caption{Sprint 3 Multi-Threading Architecture}
\end{figure}

\section{UI Updates for Multi-File Results}
The secondary view was updated to display:
\begin{itemize}
    \item Total files processed
    \item Success/failure counts
    \item Total words and characters across all files
    \item Parallel processing time
    \item Speedup factor (sequential time / parallel time)
    \item Threads used
    \item TableView with individual file results
\end{itemize}

\section{Performance Metrics}
\begin{table}[H]
\centering
\begin{tabularx}{\textwidth}{|l|X|}
\hline
\textbf{Metric} & \textbf{Description} \\
\hline
Parallel Time & Wall-clock time for concurrent execution \\
\hline
Sequential Equivalent & Sum of individual file processing times \\
\hline
Speedup Factor & Sequential Time / Parallel Time \\
\hline
Threads Used & Number of active worker threads \\
\hline
\end{tabularx}
\caption{Performance Metrics}
\end{table}

\section{Sprint 3 Deliverables}
\begin{itemize}
    \item[$\checkmark$] ThreadPoolManager with ExecutorService
    \item[$\checkmark$] MultiFileAnalysisService for concurrent processing
    \item[$\checkmark$] Thread-safe result collection with ConcurrentHashMap
    \item[$\checkmark$] Progress callbacks for real-time UI updates
    \item[$\checkmark$] Multi-file results UI with TableView
    \item[$\checkmark$] Performance metrics (speedup factor)
\end{itemize}

%========================================
% CHAPTER 5: TESTING & RESULTS
%========================================
\chapter{Testing and Results}

\section{Test Scenarios}

\subsection{Single File Analysis}
\begin{itemize}
    \item Load one text file
    \item Verify word count, character count, and frequency analysis
    \item Check progress bar animation
    \item Confirm UI responsiveness
\end{itemize}

\subsection{Multi-File Analysis}
\begin{itemize}
    \item Load 3-5 text files of varying sizes
    \item Verify concurrent processing (multiple thread names in logs)
    \item Check speedup factor $> 1.0$
    \item Confirm all files processed successfully
\end{itemize}

\subsection{Error Handling}
\begin{itemize}
    \item Attempt to load non-existent file
    \item Attempt to load non-text file
    \item Verify graceful error display
\end{itemize}

\section{Performance Results}
\begin{table}[H]
\centering
\begin{tabularx}{\textwidth}{|l|l|l|l|X|}
\hline
\textbf{Files} & \textbf{Sequential} & \textbf{Parallel} & \textbf{Speedup} & \textbf{Threads} \\
\hline
1 file & 150ms & 155ms & 0.97x & 1 \\
\hline
3 files & 450ms & 180ms & 2.50x & 3 \\
\hline
5 files & 750ms & 220ms & 3.41x & 4 \\
\hline
8 files & 1200ms & 350ms & 3.43x & 4-8 \\
\hline
\end{tabularx}
\caption{Performance Comparison (Sample Data)}
\end{table}

%========================================
% CHAPTER 6: CONCLUSION
%========================================
\chapter{Conclusion}

\section{Summary}
This project successfully implemented a Multi-Threaded Text Analyzer application through three development sprints:

\begin{enumerate}
    \item \textbf{Sprint 1} - Created a modern, user-friendly JavaFX interface with file loading capabilities
    \item \textbf{Sprint 2} - Implemented text analysis algorithms with background processing using JavaFX Task
    \item \textbf{Sprint 3} - Introduced multi-threading with ThreadPoolExecutor for concurrent file analysis
\end{enumerate}

\section{Key Learnings}
\begin{itemize}
    \item JavaFX application architecture (FXML, Controllers, CSS styling)
    \item Background thread execution to maintain UI responsiveness
    \item Thread pool management with ExecutorService
    \item Thread-safe data structures (ConcurrentHashMap, AtomicInteger)
    \item Future/Callable pattern for async result collection
\end{itemize}

\section{Future Improvements}
\begin{itemize}
    \item Sentiment analysis using NLP libraries
    \item Export results to PDF/CSV
    \item File comparison analysis
    \item Word cloud visualization
    \item Support for more file formats (PDF, DOCX)
\end{itemize}

%========================================
% APPENDIX
%========================================
\appendix
\chapter{Installation Guide}

\section{Prerequisites}
\begin{itemize}
    \item Java 21 or higher
    \item Maven 3.8+
    \item JavaFX 21 (managed by Maven)
\end{itemize}

\section{Build and Run}
\begin{lstlisting}[language=bash]
# Clone repository
git clone https://github.com/MarwaneQ/text-analyzer.git
cd text-analyzer

# Build the project
mvn clean compile

# Run the application
mvn javafx:run
\end{lstlisting}

\chapter{Module Configuration}
\begin{lstlisting}[language=Java, caption=module-info.java]
module com.textanalyzer {
    requires javafx.controls;
    requires javafx.fxml;

    opens com.textanalyzer to javafx.fxml;
    opens com.textanalyzer.controllers to javafx.fxml;

    exports com.textanalyzer;
    exports com.textanalyzer.controllers;
    exports com.textanalyzer.models;
    exports com.textanalyzer.services;
    exports com.textanalyzer.exceptions;
}
\end{lstlisting}

\end{document}
